% !TeX root = ../教材题库_全部条目_转写工作稿.tex
% Source: XJ-B1-C04-PR-P146-Y0993
% 第4章 / 4.5 函数模型及其应用 / 教材 p.146 / PDF 第153页
% OCR未校对：公式与文字均可能有误；未完成数学TeX转写；禁止直接发给学生
% 原题及相邻原解截图（核对用）：../题目原图/XJ-B1-C04-PR-P146-Y0993-01.webp
% 本文件只在转写完成后作为题干使用；OCR 条目尚不是可靠题干。
\begin{bankitem}{p.146 练习 第1题}{第4章 / 4.5 函数模型及其应用 / 教材 p.146 / PDF 第153页}{XJ-B1-C04-PR-P146-Y0993}{OCR未校对：公式与文字均可能有误；未完成数学TeX转写；禁止直接发给学生}
\begin{ocrdraft}
[教材p.146 / PDF第153页]
\par
1.从国家统计局网站可以了解到中国居民2021一2024年年末固定互联网宽
\par
带接入用户数约为下表所示：
\par
2021
\par
2022
\par
2023
\par
2024
\par
年末固定互联网宽带接入用户数/万户53579
\par
58965
\par
63631
\par
66983
\par
(1)描点画出这四年年末固定互联网宽带接入用户数随年份变化的大致
\par
图象；
\par
（2）试选择合适的函数类型建立数学模型，刻画这四年年末固定互联网宽带
\par
接入用户数随年份变化的规律，并预测2025年年末固定互联网宽带接入用户数．
\end{ocrdraft}
\end{bankitem}
